\section{Introduction}

Mobile ad hoc networks (MANETs) operate without fixed infrastructure. Nodes
organise themselves, forward one another's traffic, and reconfigure as the
topology changes. The properties that make this attractive for disaster
response, defence communication and location-aware services are the same
ones that make the network hard to secure. There is no central authority to
consult, no stable membership, and no reliable way to tell a node that is
failing from a node that is lying.

Beacon nodes deserve separate treatment within this problem. A beacon
advertises position information that neighbouring nodes consume to
localise themselves. A compromised beacon therefore corrupts the position
estimates of nodes it never exchanged data traffic with, and the corruption
propagates through the localisation computation, not through the
routing fabric. The consequences differ in kind from ordinary
packet-dropping attacks, not merely in degree. A node that drops packets
degrades a service the network can measure; a beacon that lies about its
position degrades a service the network has no independent way to check.

Elimination of malicious beacons from localisation evidence is well studied
for static beacons. Liu, Ning and Du detect a beacon whose RSSI-implied
distance disagrees with the distance implied by its declared position, and
revoke it once enough alerts accumulate~\cite{liu2005detecting}. Kuriakose
et al.~\cite{kuriakose2023embn} address the same problem in EMBN-MANET, a
UWB-based mobile ad hoc network, using RSSI and trilateration residuals, and
report complete identification when fewer than four malicious beacons share
a neighbourhood. Both establish that detection from localisation evidence is
tractable. What neither settles is the decision that precedes removal, and
what happens to the network when that decision is wrong: revocation in both
is permanent.

Both questions matter for a practical reason. In a mobile wireless
environment, congestion, transient link failure, interference and mobility
all produce observations that resemble attack. A verification stage
separating detection from confirmation therefore has a plausible role, a
need also identified in work on trust change and revision~\cite{chatzidakis2022trust}
and on verification supported by observable evidence~\cite{li2019dapv,ryu2023reputation}.
Because a removed beacon is a beacon whose localisation service is lost
to its neighbours, removal is not free: a verification mechanism that is too
eager degrades exactly the service it exists to protect.

\subsection{Problem statement}

The paper addresses four questions. The first is whether a detector that
uses one class of evidence can confirm beacons whose attack leaves no
signature in that class, and whether collecting both behavioural and
localisation evidence closes the gap (Sections~\ref{sec:coverage},
\ref{sec:baselines}, \ref{sec:sparse} and~\ref{sec:scale}). The second is how
much honest-beacon service isolation policies withdraw, and how waiting for
confirmation changes that cost (Sections~\ref{sec:policy}
and~\ref{sec:phi}). The third is whether an isolated beacon can be released
safely once the evidence against it clears (Section~\ref{sec:release}). The
fourth is when fused evidence can confirm an attacker that neither branch
confirms alone, and whether a fused detector is better than classical
single-evidence detectors run side by side (Sections~\ref{sec:region},
\ref{sec:sweep}, \ref{sec:baselines}, \ref{sec:sparse} and~\ref{sec:scale}).

\subsection{Contributions}

We answer these questions in ns-3 across 780 simulation runs: a 30-seed
detection campaign over three attack scenarios and a clean network, six
isolation policies replayed over that campaign and checked by enforcing them
in the simulator, an iterated closed loop, two attack-intensity sweeps,
log-normal shadowing, a sparse topology, a 50-node network with two
simultaneous attackers, and two re-implemented literature baselines evaluated
on the same traces.

\begin{enumerate}
\item \textbf{Evidence coverage.} We frame malicious-beacon
      verification as a question of evidence coverage rather than detection
      accuracy. Every detector restricted to one evidence class, including
      re-implemented localisation-only and behaviour-only baselines from the
      literature, was blind to the attack outside that class, in the small,
      sparse and 50-node networks alike; collecting both classes confirmed
      the attacker in 209 of 210 attacked runs of the dense network and
      both attackers in every run at 50 nodes. The verification pipeline that
      fuses the two classes into per-observer trust and a three-state verdict
      is built from established components and is the vehicle for
      the study.

\item \textbf{A measure of what false isolation costs.} We introduce
      cumulative false isolation~$\Phi$, in node-seconds, and its normalised
      form~$\bar\Phi$, as complements to the false-positive rate, which says
      nothing about how long an error persists. Isolating only on
      confirmation withdrew 3.9--5.0\% of honest-beacon service on a clean
      network against 74.8\% for the strongest immediate-isolation policy,
      and the gap held when the schedules were enforced in the simulator,
      under shadowing and at 50 nodes.

\item \textbf{Release is as consequential as isolation.} Timed
      redemption is known to let attackers back in~\cite{bansal2003ocean,sen2010survey}.
      We show that release is unsafe even without a timer or any action by
      the attacker: releasing an isolated beacon when its verdict clears
      re-admits forwarding attackers in a closed loop, because isolation
      removes the very traffic their behavioural evidence comes from. An evidence-aware release rule
      prevented every re-admission at a small cost in false isolation, and
      showed that the advantage of confirmation-gated isolation comes from
      the gating itself.

\item \textbf{When fusion helps.} Every weighted Stouffer
      detector has a fusion-only region, containing attackers whose two anomalies are
      each sub-threshold and of comparable size, and an attack family scaled
      in intensity passes through it only if the branches' detection
      thresholds for that family lie within a factor of
      $1+\sqrt{2}\approx 2.414$. We offer this as an explanatory tool rather
      than a new result in fusion theory. The condition is checkable before any
      fusion experiment and predicted the outcome of two intensity sweeps.
      We also make explicit, for trust fusion, the calibration requirement on
      which any such test depends~\cite{stouffer1949,whitlock2005weighted}.

\item \textbf{Comparison with classical detectors.} In the
      small dense network a calibrated disjunction of the two literature
      baselines matched the fused pipeline's coverage and false isolation
      and confirmed attackers as fast or faster, as the fusion condition
      predicts. In the sparse and 50-node networks the same disjunction
      withdrew 11--24\% of honest-beacon service on a clean network against
      0--3\% for the pipeline. The pipeline's single-branch results show
      why: in the sparse network its peer-relative scoring of behaviour
      avoided the false alarms, and at 50 nodes fusion suppressed false
      alarms raised by its own behavioural branch. Fusion thus trades
      sensitivity near the threshold for fewer false isolations. In the
      sparse network the pipeline paid for this by confirming the
      forwarding attacker in only 4 of 10 seeds.

\item \textbf{Bounding results.} Localisation error
      (Appendix~\ref{sec:locerr}) shows that the two
      attack classes are disjoint in both evidence directions; detection and
      the false-isolation result survive 4 and 8\,dB shadowing at a cost in
      latency for forwarding attacks; and we report two negative results:
      the trust aggregation stage changed no verdict
      (Appendix~\ref{sec:ablation}), and neither isolation policy produced a
      throughput advantage surviving its confidence interval
      (Section~\ref{sec:policy}).
\end{enumerate}

\subsection{Organisation}

Section~\ref{sec:related} positions the work against the closest prior art.
Section~\ref{sec:model} defines the network and threat models.
Section~\ref{sec:framework} develops the model in full.
Section~\ref{sec:method} describes the simulation campaign.
Section~\ref{sec:results} reports results, Section~\ref{sec:discussion}
interprets them, Section~\ref{sec:threats} states threats to validity, and
Section~\ref{sec:conclusion} concludes. Appendix~\ref{sec:appendix}
reports secondary results.
